\chapterNo{Résumé}
\addcontentsline{toc}{chapter}{Résumé}
\medskip	

Ce mémoire a été rédigé dans le cadre du Master Technologies numériques appliquées à l'histoire de l'École nationale des chartes en vue de la validation de ce diplôme. Il a été réalisé suite à un stage de quatre mois au sein du projet \gls{eida} à l'Observatoire de Paris qui étudie les diagrammes astronomiques de tradition ptoléméenne du \VIII au \XVIII siècle avec l'utilisation de la vision artificielle. Ce mémoire présente le développement de preuves de concept pour la création de visualisations afin de faciliter l'analyse des résultats des traitements automatiques. \\

This thesis was written as part of the Master's program in Digital Technologies Applied to History at the École nationale des chartes, in order to fulfill the requirements for this degree. It was completed following a four-month internship within the \gls{eida} project at the Observatoire de Paris, which studies Ptolemaic astronomical diagrams from the 8th to the 18th century using computer vision. This thesis presents the development of proof-of-concept implementations for the creation of visualizations aimed at facilitating the analysis of results from automated processing.\\

\textbf{Mots-clés~:} histoire des sciences; astronomie; visualisation; réseaux; vision artificielle; Gephi; D3.js.\\

\textbf{Informations bibliographiques~:} Lucie Ledieu, \textit{Explorer les réseaux de transmission de données historiques : Elaboration de visualisations pour l'application AIKON}, mémoire de master \og Technologies numériques appliquées à l'histoire~\fg, dir. Ségolène Albouy, École nationale des chartes, 2025.

\clearemptydoublepage